% =====================================================================
% sections/boundary.tex — §5.8 泛化边界（rev.1，2026-07-07 轻流程试点落地）
% 依据：spec §2 §5.8 块 + next_session_prompt.md（轻流程：论点级骨架 → 直接落地）。
% 数字 ground truth（唯一权威）：docs/rebrac_broad_validation_v2_report.md
%   §2（N0 0.850±0.024 / per-seed 0.867, 0.833 / Δ−5.2pp vs 0.902±0.021）
%   §3（N2' 0.000/0.000, 0/60, OOB 59/60, progress≈0.018, 航程≈88 m；
%       crosscomp 临界 sanity 0/30, progress −0.682, 超时 8）
%   §4.5（asym ablation 0/60；配对 mean Δ+0.1025=1.15 SEM / median −0.016 /
%       单回合 46% / 17/30 |Δ|<0.10；verdict ACTOR_FUNDAMENTAL_CONFIRMED）
%   §5（上限分解、oracle direct 0.700=21/30、70pp 差距、claim 边界）。
%   在线参照 0.26±0.15（k=4）/ 0.88±0.04（k=12）取自 sections/online.tex
%   §5.5.3 瓶颈表（spec 缝合点：对照须用 0.26，非 legacy 10% 单种子值）。
% 先决裁决（2026-07-07 用户拍板）：① N2' n=2 以极重 caveat 呈现（显式 n=2 +
%   rule-of-three 上界 + 补种子登记语，见 §5.8.5），补种子不阻塞写作；
%   ② 在线特权消融不与本节并列，仅散文指回 §5.5（线别限定、不下跨线结论），
%   特权 critic 四态统一收束留 §5.10。
% 红线（spec §2 §5.8 + §0.4）：✅「特权价值网络不能挽救 / 排除价值侧成因 /
%   加固策略侧解释」；❌「信息论不可能 / proven actor-incapable」（零成功兼容
%   表征不足与机制未转化两读法，§5.8.5 明写）；本节发现不嫁接回 §5.7；负面
%   结果作 finding 同台呈现统计强度（n / 预登记判定 / 零计数上界）；生物不
%   回扣、零跨章引述；§5.8.m 只写本节增量，协议指针回 §5.3/§5.4。
%
% rev.2（2026-07-07）：独立对抗复审落地（轻流程试点检验轮，对抗+系统性合并；
%   findings 存档 section_5_8_review_findings.md）。数字忠实性零漂移（全量回查
%   report §2–§5/§4.5 + plan 预登记原文 + online/rebrac/setup 刊值，逐项一致）；
%   红线零踩线。修订实质项（0 CRIT / 1 HIGH / 5 MED）：
%   [HIGH-披露] plan §6.2 预登记三种子、首轮执行两种子（report 头部 seed-count
%     caveat），成稿以「预登记」背书却未披露缺口 → §5.8.m 补如实登记（缺额并入
%     定稿前补充验证事项）；§5.8.5「预登记门槛据此不要求追加种子」收窄归属为
%     「消融的预登记判读亦未触发补种子复核」（该混合区触发条款属消融 gate，
%     不为 N0/N2' 的种子缩水作豁免）。
%   [MED-端点] §5.8.3 消融门槛 (0.10,0.40] / ≥0.40 端点双重归属且与 §5.8.m 矛盾
%     → 统一为 (0.10,0.40) 混合、0.40 及以上重构（report 原文端点粗糙，取自洽读法；
%     结果 0.000 距端点极远，不影响判定）。
%   [MED-衔接] §5.8.m「相应工况固定评估集的全部 30 回合」与 §5.3.6「主工况主评估
%     集 100 回合」相抵 → 改写：两单元各用其工况下 30 回合固定评估集，亚临界侧与
%     主线主评估集同任务分布、独立采样、规模不同（30 对 100）、样本不重合。
%   [MED-全称] §5.8.2「不存在直接执行合格的可部署规则采集器」「本章可用采集器的
%     最高水平」超出实测（goalseek/worldcomp 临界未实测，plan §8 仅登记子集逻辑
%     豁免）→ 改为显式推理承载（较强的横流补偿参照已零成功、目标直指参照无补偿
%     能力）+「实测采集器中的最高水平」。
%   [MED-限定] §5.8.5 末段「完整轮廓」→「沿流动工况轴勾勒」+ 补横流穿越任务与
%     k=4 观测构成限定（与行 97 自划边界及 §5.7.5 三重限定对齐）。
%   [MED-语体] 开篇「两个边界…一是…二是…」枚举（§0.5.9(a) 结构暴露）→ 散文对举。
%   LOW×6 仅记录（表内简称不一/「在线 SAC」行名/对照口径列举/≥0.40 强正面区间
%     省略/数据集行指称歧义/「定位在策略一侧」措辞），见 findings §4，收尾统稿轮处理。
% =====================================================================

\section{泛化边界：临界工况下可部署性能的策略侧上限}\label{sec:ch5_boundary}

\S\ref{subsec:ch5_rebrac_implications} 把离线主线的结论界定在亚临界工况、历史效率导向奖励与横流穿越几何之内，并把可部署性能在更强流动工况下能否保持留作本节的问题。本节把离线主线方法推到其适用范围的边界上考察：一个边界沿奖励设定展开——由历史效率导向换为到达导向后，主线结论是否在亚临界工况保持；另一个边界沿流动工况展开——由亚临界（$\mathrm{Re}=150$、$U_\infty=1.0$ m/s）加深到临界（$\mathrm{Re}=250$、$U_\infty=1.5$ m/s，速度比 $\lambda=1.0$）后，以可获得的最高品质示范数据训练，可部署单点观测策略能否继续工作。前者为后者提供锚点——若奖励迁移本身即破坏主线性能，临界工况的失效便无从归因；后者给出全章最强的负面发现，其失效归属再由一项单变量消融判别。

两个考察单元沿用 \S\ref{subsec:ch5_method_rebrac} 的 ReBRAC-Q 与主线锁定的行为约束配置 $(\beta_1,\beta_2)=(4.0,\,2.0)$，各以两个随机种子训练，在相应工况的固定评估集上以确定性策略终检评估（每种子 $30$ 回合，\S\ref{subsec:ch5_eval}）；全部判定门槛在实验执行前预先登记，结果按门槛判读而非事后解释。本节种子数为二、少于主线的五，而其临界结论又是全章分量最重的负面结论；因此除以预登记门槛与零计数上界呈现统计强度外，本节全部结论的证据分量都应按两种子解读，\S\ref{subsec:ch5_boundary_implications} 对此专门讨论。表~\ref{tab:ch5_boundary_cells} 汇总两单元及消融的逐种子结果与判定。

\begin{table}[t]
\centering
\caption{泛化边界考察的两个单元与机制判别消融。每单元两个随机种子、每种子终检 $30$ 回合，判定门槛均为实验前预登记；两单元同用 ReBRAC-Q 完整配置 $(\beta_1,\beta_2)=(4.0,\,2.0)$ 与到达导向奖励，锚点单元在亚临界评估集、其余两行在临界评估集上评估；消融的唯一变量为价值网络按特权协议读取等效来流。零成功行按 \S\ref{subsec:ch5_eval} 的零计数口径给出 $95\%$ 置信上界。}
\label{tab:ch5_boundary_cells}
\footnotesize
\setlength{\tabcolsep}{5pt}
\begin{tabular}{l l c c l}
\toprule
\textbf{单元} & \textbf{训练数据} & \textbf{种子 42 / 0} & \textbf{均值} & \textbf{预登记判定} \\
\midrule
亚临界锚点 & 横流补偿参照，亚临界 & $0.867$ / $0.833$ & $0.850 \pm 0.024$ & $\ge 0.70$：保持 \\
临界单元 & 全场走廊参照，临界 & $0.000$ / $0.000$ & $0.000$（上界 $0.05$） & $< 0.15$：强负面 \\
临界 $+$ 特权价值网络 & 同上（价值网络读等效来流） & $0.000$ / $0.000$ & $0.000$（上界 $0.05$） & $\le 0.10$：排除价值侧 \\
\bottomrule
\end{tabular}
\end{table}

\subsection{亚临界锚点：到达导向奖励下主线结论的保持}\label{subsec:ch5_boundary_anchor}

到达导向奖励下，以横流补偿参照的一千回合亚临界数据（\S\ref{subsec:ch5_datasets}）训练的 ReBRAC-Q 可部署策略，在亚临界横流穿越评估集上的成功率为 $0.850 \pm 0.024$（两种子分别为 $0.867$ 与 $0.833$），高于预登记的保持门槛 $0.70$。主线的可部署离线性能没有因奖励设定更换而崩解，后续临界工况的失效因此不能归因于奖励迁移本身。

与主线在历史效率导向奖励下的 $0.902 \pm 0.021$（\S\ref{subsec:ch5_rebrac_scale}，五种子）相比，锚点单元低约 $5.2$ 个百分点，且两种子同向退化（分别约 $-3.5$ 与 $-6.9$ 个百分点）。同向性提示这更可能是到达导向塑形下策略略趋激进、越界回合略增的系统性效应，而非种子噪声；但两种子的同向本身证据力有限，且该对照跨越奖励设定、种子数与终检回合数三个口径（五种子对两种子、每种子 $100$ 回合对 $30$ 回合），只宜作量级参考。对本节而言，锚点的作用不在精确定标退化幅度，而在确认到达导向奖励下主线方法仍处于正常工作区。

\subsection{临界工况：高品质示范数据不能支撑可部署策略}\label{subsec:ch5_boundary_critical}

临界工况下，两个严格单点观测的规则采集器中较强的横流补偿参照在临界评估集上零成功（$0/30$，超时 $8$、越界 $22$），其固定规则的偏航补偿在临界涡街的扰动下无法完成闭环校正；更弱的目标直指参照不具任何补偿能力，可部署规则采集器由此整体出局。临界单元的离线数据因此由全场走廊参照采集——该控制器读取全场流场为候选走廊打分，并以等效来流作局部来流补偿（补偿项读取的正是 \eqref{eq:ch5_priv_obs} 所定义的等效来流，仅坐标系不同），在同一评估集上直接执行成功率 $0.700$（$21/30$），是该工况下本章实测采集器中的最高水平。这也意味着两单元之间不构成单变量对照——工况与采集器同时改变——本节的推断不依赖两单元相减，而是各自与其预登记门槛比较。

以该示范数据（一千回合）训练的 ReBRAC-Q 可部署策略在临界评估集上零成功：两种子合计 $0/60$，按 \S\ref{subsec:ch5_eval} 的零计数口径，$95\%$ 置信上界约 $0.05$，落入预登记的强负面区（$<0.15$），对采集器直接执行水平的恢复比例为零。两种子的零成功完全一致，不属于评估噪声边缘的低值，而是完全失效。

失效形态本身携带机制信息。六十个回合中五十九个以越界终止、仅一个超时；策略平均航程约 $88$ m，回合级净接近比例（初末目标距离之差与初始距离之比）均值仅约 $0.02$，且回合间离散极大（种子内标准差约 $0.5$ 至 $0.6$）。与横流补偿参照的失效形态（净接近比例 $-0.68$、超时占 $8/30$，近乎在原地被涡街推离）对照，学得的策略显然习得了示范数据中的运动模式——它持续游动、大致朝向目标——但缺少把航迹收敛到目标的闭环校正，轨迹在扰动下振荡发散直至越界：这是闭环失稳的足迹，而非策略停滞的足迹。

与全部观测一致的解释指向示范数据的信息结构。全场走廊参照的闭环校正以等效来流为条件，而临界工况下单点瞬时采样与等效来流的对应关系微弱——\S\ref{subsec:ch5_online_ablation} 的在线单变量证据独立支持这一点：$k=4$ 时序窗下的单点观测不足以支撑闭环校正。于是以可部署观测为条件的行为克隆项~\eqref{eq:ch5_rebrac_actor} 所能拟合的，至多是示范动作在该条件下的平均模式：同一个可部署观测对应示范者多个不同的校正动作，条件平均把闭环校正抹平成大致朝目标游动的开环倾向。失效并非因为示范不好，而是因为示范的决策依据落在策略观测不到的维度上。

\subsection{机制判别：特权观测注入价值网络不能挽救失效}\label{subsec:ch5_boundary_asym}

零成功本身无法区分两类成因：或在策略侧——可部署观测下策略无从表达闭环校正；或在价值侧——价值网络同样只读可部署观测，其价值估计在临界工况失真，令自举目标劣化、训练本身失败。为隔离二者，消融在唯一变量下重训临界单元：价值网络及其目标网络按特权协议读取等效来流（\S\ref{subsec:ch5_method_protocol}），策略网络保持单点观测、其改进步骤中特权通道置零以对齐部署接口；数据集、行为约束配置、种子组与训练预算与临界单元逐项相同，终检使用同一评估集与同一评估随机序列，两组结果因此可逐回合配对。

结果没有任何改变：两种子仍为 $0/60$（零计数上界同前，约 $0.05$），终止构成仍由越界主导（一个种子三十回合全部越界，另一种子二十九个越界加一个超时）。按预登记门槛——成功率 $\le 0.10$ 判排除价值侧成因，$(0.10,\,0.40)$ 判混合、须补种子复核，$0.40$ 及以上判重构解释——结果落入第一区间。

逐回合配对进一步排除成功率看不出、行为上却有改善的余地。一个种子的净接近比例均值由 $0.041$ 升至 $0.143$，表面似有提升；但配对差的均值 $+0.10$ 仅合 $1.15$ 倍标准误，配对差的中位数为 $-0.016$，均值差的近半（约 $46\%$）来自单一回合（该回合在基线下超时跑飞、消融下普通越界），三十个配对回合中十七个的差异绝对值不足 $0.10$，两种子的终止构成差异方向还彼此相反。给价值网络以完整的等效来流，可部署策略的部署行为与基线在统计上不可区分。

这一判别把临界失效的归属从价值侧移开：瓶颈不在价值估计的质量——训练全程向价值网络注入部署时不可得的等效来流，也没有把策略抬离完全失效的水平——失效与策略侧解释一致，并使该解释得到加固。

\subsection{性能上限的分解与在线参照对照}\label{subsec:ch5_boundary_ceiling}

表~\ref{tab:ch5_boundary_ceiling} 把临界工况同一评估集上的各层水平并置。可部署接口一侧，规则控制器与离线强化学习——无论价值网络读取什么——均为零，在线强化学习约 $0.26$；不可部署一侧，全场走廊参照直接执行 $0.700$。采集器与以其示范训练的可部署策略之间 $70$ 个百分点的差距，量化了临界工况、单点观测下本任务的部分可观测差距；本节的消融证据把这一差距定位在策略一侧。

\begin{table}[t]
\centering
\caption{临界工况固定评估集上的性能层级分解。可部署各行的策略在部署时仅读单点观测；各行训练协议互不相同——在线行取自 \S\ref{sec:ch5_online} 的瓶颈分析（三种子、每次交互 $100$ 万步），离线两行为本节单元（两种子），两个参照策略为确定性规则控制器单次评估——本表刻画层级而非受控比较。}
\label{tab:ch5_boundary_ceiling}
\small
\begin{tabular}{l l c}
\toprule
\textbf{层} & \textbf{训练期信息来源} & \textbf{成功率} \\
\midrule
横流补偿参照（直接执行） & 无学习；单点流速规则补偿 & $0.000$（$0/30$） \\
ReBRAC-Q，可部署协议（两种子） & 全场走廊参照示范一千回合 & $0.000$（$0/60$） \\
ReBRAC-Q，特权协议价值网络（两种子） & 同上 $+$ 价值网络读等效来流 & $0.000$（$0/60$） \\
在线 SAC，$k=4$（三种子） & 在线交互一百万步，单点观测 & $0.26 \pm 0.15$ \\
全场走廊参照（直接执行，不可部署） & 全场流场 $+$ 等效来流 & $0.700$（$21/30$） \\
\bottomrule
\end{tabular}
\end{table}

表中的次序本身构成一个发现：以可获得的最高品质示范做离线训练，反而低于同一评估集上以相同单点观测与 $k=4$ 历史直接在线交互的标准 SAC（$0.26 \pm 0.15$，\S\ref{subsec:ch5_online_ablation}）。与失效形态相容的一个解释是，在线策略在其自身经历的状态分布上学习，行为与其可观测信息自洽；离线的行为克隆项则把策略锚定到以不可观测量为条件的示范动作分布上，评估中一旦偏离示范航迹，策略便被拉向其未见过的状态区域上的动作先验，行为发散。该解释与上一小节的证据一致，但这一对照跨线而非受控——两侧的训练预算、种子组与实施细节均不同——只作方向性读数，不作效应量比较。

在线一侧的证据同时划定了这一边界的适用范围。同一临界工况下，把观测历史延长到 $k=12$ 可使在线成功率升至 $0.88 \pm 0.04$（\S\ref{subsec:ch5_online_ablation}）：单点观测在时间维度上携带着足以闭合差距的信息。但本节离线数据以 $k=4$ 观测构成采集，离线学习只能复用既采数据的观测构成，在线可行的时序救援途径在本节设定内不可用。因此本节所划的是 $k=4$ 观测构成下、临界工况离线学习的经验上限，不外推到更长时序窗的离线数据；把时序利用引入离线一侧属于策略侧改造方向，超出本节闭环。

\subsection{对中心命题的含义与结论边界}\label{subsec:ch5_boundary_implications}

本节给出中心命题反向一侧分量最重的负面发现。三种为系统增添能力的手段中，向价值网络注入特权观测在此触底：亚临界主线上它已退为非必需（\S\ref{subsec:ch5_rebrac_implications}），临界工况上它连挽救完全失效的能力也没有——价值侧的信息注入没有转化为可部署策略的性能。在线一侧的对应消融在其自身协议内同样未闭合临界差距（\S\ref{subsec:ch5_online_ablation}）；特权信息在两条线索、多种工况下呈现的不同失效方式，其统一收束留待本章末节。

结论的边界须精确划定。本节证据支持的表述是：接收特权观测的价值网络不能挽救临界工况下可部署策略的失效，价值侧成因被排除，策略侧上限的解释被加固。本节证据不支持、本章也不采用的表述是：单点观测策略被证明不可能学得该任务。消融的零成功同时兼容两种读法——$k=4$ 观测构成下的单点观测在表征上不足以承载闭环校正；或非对称价值网络这一机制本身未把特权信息有效转化到策略改进中（策略改进时特权通道置零，价值网络的信息优势只能间接经由动作评价传导）。区分二者需要策略侧的受控改造，如更长时序窗、循环结构或显式的流场估计通道，不在本节范围之内。

种子数是本节结论最须直面的限制。两个单元与消融各以两个随机种子完成，是全章种子数最少的一组实验，而临界结论又是全章分量最重的负面结论。支持在两种子下报告的理由在于信号形态：两种子给出完全一致的退化零成功（合计 $0/60$、零计数上界 $0.05$），是零方差的强信号而非检验力不足的弱信号，消融的预登记判读亦未触发补种子复核（触发条件为落入混合区）。尽管如此，本节结论的证据分量仍应严格按两种子解读；临界单元与消融在更多种子下的配对复核已列为本章定稿前的补充验证事项，在其完成之前，引用本节临界结论时应连同两种子这一限定一并引用。

本节的边界发现不改写前两节的结论。亚临界工况、历史效率导向奖励下的主线发现由其自身协议支持，本节锚点单元又确认了其在到达导向奖励下的保持；反过来，主线的正面结果也不为临界工况提供任何担保。两端合起来沿流动工况轴勾勒出可部署离线策略的适用范围：横流穿越任务与 $k=4$ 观测构成下，亚临界工况内成立，临界工况处触及策略侧的经验上限。

\subsection{实施细节与可复现性}\label{subsec:ch5_boundary_impl}

本节两个训练数据集为表~\ref{tab:ch5_datasets} 中到达导向奖励下的亚临界与临界两行：亚临界集由横流补偿参照采集、临界集由全场走廊参照采集，各一千回合；两集均为 $s_0$ 传感、$k=4$ 历史堆叠（观测 $48$ 维），并随转移记录特权观测及其后继（\S\ref{subsec:ch5_impl}）。两个采集器在临界评估集上的直接执行水平（全场走廊参照 $21/30$、横流补偿参照 $0/30$）以与终检相同的固定评估集和确定性执行测得。

训练配置与主线一致：超参数取表~\ref{tab:ch5_rebrac_hparams} 的完整配置（含价值网络层归一化），每种子训练 $64$ 个训练轮次。与主线正式协议的差异有三处：种子组为两种子而非五种子——本节所依计划预登记为三种子，首轮实际执行为两种子，缺额的补足已并入 \S\ref{subsec:ch5_boundary_implications} 登记的定稿前补充验证事项；不作验证集字典序检查点选择，直接评估训练终点检查点——本节考察失效边界而非最优配置，终点口径对各单元一视同仁；终检各在相应工况的 $30$ 回合固定评估集上进行，其中亚临界一侧的评估集与主线的每种子 $100$ 回合主评估集取自同一任务分布、但为独立采样且规模不同，样本并不重合，\S\ref{subsec:ch5_boundary_anchor} 的跨口径对照因此只作量级参考。锚点单元与临界单元除数据集与评估集外配置完全相同。

非对称消融沿 \S\ref{subsec:ch5_method_protocol} 的特权协议实现：价值网络及其目标网络的输入扩展为 $(\mathbf{o},\mathbf{a},\mathbf{o}^{\mathrm{priv}})$，训练全程读取数据集记录的特权观测及其后继；策略改进步骤中特权通道以零向量填充，使策略梯度所经的价值面与部署一致。除该单一变量外，数据集、超参数、种子组与训练轮次与临界单元逐项相同；终检使用同一评估集与同一评估随机序列，逐回合配对由此成立，配对诊断中的标准误按种子内配对差的回合样本计算。

全部判定门槛于实验执行前登记：锚点单元以 $0.70$ 为保持门槛；临界单元以 $0.15$ 为强负面上界，$[0.15,\,0.40]$ 为触发行为约束系数扫描的中间区（实际结果未触发）；消融以 $0.10$ 为排除价值侧门槛，$(0.10,\,0.40)$ 为须补种子复核的混合区（未触发），$0.40$ 及以上为重构解释区。零计数上界按 \S\ref{subsec:ch5_eval} 的口径以 $3/n$ 计（$n=60$ 时约 $0.05$）。在线参照数字取自 \S\ref{sec:ch5_online} 的临界瓶颈分析（三种子终检的 $k=4$ 与 $k=12$ 配置），其与本节协议的差异已随正文标明。
